
\section{Introduction}

\subsection{Background and Motivation}

As Wi-Fi infrastructure has become nearly ubiquitous in modern indoor environments, it has evolved from a communication utility into a versatile medium for passive sensing. This capability enables advanced applications such as human activity recognition (HAR) and gait-based individual identification, offering a more secure and privacy-preserving alternative to conventional biometric systems that can be easily spoofed or visually manipulated \parencite{shi2021wifi}. 

The foundation of these capabilities lies in the unique properties of Wi-Fi Channel State Information (CSI), which encodes minute perturbations in the wireless channel caused by human motion. These subtle variations are extremely difficult to replicate or forge, rendering CSI-based biometric identification inherently robust against impersonation attacks \parencite{alsheikh2024wifiid}. Moreover, because CSI captures fine-grained dynamics related to body posture, stride, and micro-movements, it exhibits strong discriminability across individuals, enabling precise gait modelling \parencite{zhang2020gaitid}. 

Importantly, Wi-Fi–based sensing is passive, device-free, and non-intrusive, making it suitable for environments where privacy, stealth, or minimal hardware deployment are critical. Applications extend across security monitoring, smart ambient environments, health-care analytics, and continuous authentication \parencite{lv2017wii}. The fact that these systems operate using existing Wi-Fi routers further enhances scalability and cost-efficiency.

\subsection{Research Context and Dataset Overview}

This thesis investigates a deep learning–based framework for individual identification through gait analysis using CSI data. The evaluation focuses on the comparative performance of four state-of-the-art architectures: CSILSTMNet, DenseNet1D, EfficientNet1D-LSTM, and MobileNetV3\_1D\_LSTM \parencite{deng2023sensefi}. 

The research leverages a dataset containing recordings from 30 participants performing 12 distinct activities. For the experiments conducted in this work, Experiment\,\#3—comprising exclusively gait-related CSI sequences totalling 1.94\,GB—was selected due to its focus on walking patterns essential for biometric identification \parencite{alsaify2020losnlos}.

\subsection{Role of MLOps and Experimentation Strategy}

A central emphasis of this study is the integration of MLOps practices to ensure reproducibility, standardized experimentation, and efficient comparative evaluation. Using MLflow, all model training runs, hyperparameter combinations, performance metrics, and learned artifacts are tracked systematically, enabling robust cross-architecture comparisons and facilitating transparent model selection \parencite{castiglioni2023whofi}. This pipeline allows the identification of the most effective framework based not only on accuracy but also on training stability, generalization capability, and computational efficiency.

Prior work on CSI-based HAR typically reports accuracies between 80--95\% when evaluated under controlled lab settings \parencite{zhuravchak2022wifi_har}. This thesis extends these investigations by evaluating performance across a broader and more diverse activity set and participant group, thereby providing insights more reflective of real-world deployment scenarios \parencite{varga2024mitigating}. 

\subsection{Technical Rationale and Deep Learning Perspective}

The proposed methodology capitalizes on the rich amplitude and phase representations embedded within CSI streams. These features encode gait-specific biometric signatures, representing a substantial advancement beyond binary presence detection or simple classification frameworks \parencite{avola2025transformercsi}. Deep learning architectures are particularly well-suited for this task because they automatically learn hierarchical representations of CSI signals, capturing nonlinear relationships that correlate strongly with human-specific motion traits \parencite{wang2019csiid}. 

By systematically comparing architectures of varying depth, width, and hybrid composition, this research addresses a critical gap in existing literature: the lack of rigorous architectural analysis in CSI-based biometric identification. Earlier studies demonstrated the feasibility of applying deep learning to CSI signals but seldom evaluated how architectural variations influence the extraction of gait-specific features or affect real-time deployability \parencite{wang2019csiid}. This thesis directly targets this research void.

\subsection{Critical Analysis of Existing Limitations}

While CSI-based sensing has demonstrated tremendous potential, the field still faces several constraints that motivate deeper investigation:

\begin{itemize}
    \item \textbf{Limited Architectural Comparisons:} Studies often focus on a single architecture without evaluating how alternative designs may better capture CSI dynamics.
    \item \textbf{Lack of MLOps Integration:} Many prior works lack systematic experimentation pipelines, making reproducibility and hyperparameter analysis difficult.
    \item \textbf{Environmental Constraints:} CSI signals are highly sensitive to router placement, wall materials, multipath effects, and human proximity, yet many works overlook these deployment-specific variations.
    \item \textbf{Dataset Diversity:} Several datasets are collected in narrow or controlled spaces, limiting generalizability across different building layouts.
    \item \textbf{Computational Inefficiency:} Complex deep learning models may achieve high accuracy but cannot be deployed on edge or real-time platforms.
\end{itemize}

These gaps underscore the need for a comparative, MLOps-driven exploration of multiple architectures under consistent settings—precisely the approach adopted in this thesis.

\subsection{Contributions of This Thesis}

The key contributions of this work can be summarized as follows:

\begin{enumerate}
    \item A systematic comparative evaluation of four deep learning models for CSI-based gait identification using a multi-participant dataset.
    \item An MLOps-driven experimental workflow leveraging MLflow for reproducibility, versioning, and transparent model comparison.
    \item Detailed analysis of model performance with respect to architectural depth, hybrid design elements, and CSI signal characteristics.
    \item Insights into the suitability of different architectures for real-time and resource-constrained deployments.
    \item A refined, deployable approach for privacy-preserving biometric authentication using everyday Wi-Fi infrastructure.
\end{enumerate}

\subsection{Chapter Summary}

This chapter introduced the core motivations behind leveraging Wi-Fi CSI for gait-based biometric identification, highlighting the limitations of traditional biometrics and the advantages of CSI as a passive and privacy-aware modality. It outlined the dataset, architectural focus, and importance of integrating MLOps for systematic experimentation. Additionally, it provided a critical assessment of existing research gaps—particularly the lack of architectural comparisons, limited dataset diversity, and insufficient reproducibility—which directly motivate the objectives of this thesis. The next chapter expands on these foundations by presenting a comprehensive review of related literature across CSI sensing, deep learning architectures, signal processing techniques, and biometric identification frameworks.




%%%%%%%%%%%%%

\section{Dataset Description}

\subsection{Overview}
This study utilizes a publicly available Wi-Fi Channel State Information (CSI) dataset designed for both Human Activity Recognition (HAR) and Human Identification (Human-ID). The dataset  comprises CSI and Received Signal Strength Indicator (RSSI) measurements collected from \textbf{30 subjects} performing \textbf{12 distinct daily activities} in controlled indoor environments. Each CSI capture provides fine-grained amplitude and phase information across multiple subcarriers, enabling the extraction of biometric gait signatures and motion-related patterns essential for classification tasks\parencite{alsaify2020losnlos}.





\subsection{Data Collection Setup}
The data was recorded using a standard Wi-Fi transmitter--receiver pair configured to extract CSI at the subcarrier level. The experimental setup includes both:
\begin{itemize}
    \item \textbf{Line-of-Sight (LOS)}
    \item \textbf{Non-Line-of-Sight (NLOS)}
\end{itemize}
conditions to simulate realistic indoor environments. CSI measurements were captured using a Wi-Fi Network Interface Card (NIC), such as the Intel 5300, capable of logging per-packet CSI information. Each CSI frame contains amplitude and phase values across multiple antenna pairs and subcarriers, along with RSSI and timestamp metadata\parencite{alsaify2020losnlos}.

\subsection{Dataset Structure}
The dataset is organized by subject and activity, where each recording session is stored as a CSV file containing:
\begin{itemize}
    \item CSI amplitude values
    \item CSI phase values
    \item RSSI measurements
    \item Timestamps
    \item Metadata fields (Activity ID, Subject ID, LOS/NLOS)
\end{itemize}

Each file consists of thousands of CSI packets, typically with:
\begin{itemize}
    \item 30--60 subcarriers per packet
    \item Multiple transmit--receive antenna combinations
    \item High-frequency temporal sampling suitable for deep learning-based time-series modeling
\end{itemize}

This rich structure enables the use of window-based segmentation (e.g., 64 or 128 samples per window with overlap), which is essential for training deep learning models.



\subsection{Activities Covered}
The dataset includes recordings of \textbf{12 daily activities}, encompassing a mixture of static, dynamic, and transitional movements. This diversity supports a comprehensive evaluation of model robustness across multiple behavioral patterns, as detailed in Table \ref{tab:Table1}.

This thesis specifically focuses on Experiment 3 (highlighted with a yellow background in the table) for capturing the gait patterns of participants.
\begin{table}[H]
\scriptsize
\centering
\caption{Experiments and activities indicators in captured Wi-Fi packets for the dataset}
\label{tab:Table1}
\begin{adjustbox}{width=1\textwidth}
    \begin{tabular}{ l l l l  }
    \hline
        \textbf{Experiment indicator} & \textbf{Experiment} & \textbf{Activity indicator} & \textbf{Description}  \\ 
        \hline
        \textbf{C1} &  \text{Falling from sitting position} & \text{A1} & \text{Sit still on a chair} \\
        \textbf{} &  \text{} & \text{A2} & \text{Falling down} \\
        \textbf{} &  \text{} & \text{A3} & \text{Lie down} \\
        \hline
        \textbf{C2} &  \text{Falling from standing position} & \text{A4} & \text{Stand still} \\
        \textbf{} &  \text{} & \text{A5} & \text{Falling down} \\
        \textbf{} &  \text{} & \text{A3} & \text{Lie down} \\
        \hline
        \cellcolor{yellow}\textbf{C3} &  \cellcolor{yellow}\text{Walking} & \cellcolor{yellow}\text{A6} & \cellcolor{yellow}\text{Walking from transmitter to receiver} \\
        \cellcolor{yellow}\textbf{} &  \cellcolor{yellow}\text{} & \cellcolor{yellow}\text{A7} & \cellcolor{yellow}\text{Turning} \\
        \cellcolor{yellow}\textbf{} &  \cellcolor{yellow}\text{} & \cellcolor{yellow}\text{A8} & \cellcolor{yellow}\text{Walking from receiver to transmitter} \\
        \cellcolor{yellow}\textbf{} &  \cellcolor{yellow}\text{} & \cellcolor{yellow}\text{A9} & \cellcolor{yellow}\text{Turning} \\
        \hline
        \textbf{C4} &  \text{Sit down and stand up} & \text{A1} & \text{Sit still on a chair} \\
        \textbf{} &  \text{} & \text{A10} & \text{Standing up} \\
        \textbf{} &  \text{} & \text{A4} & \text{Stand still} \\
        \textbf{} &  \text{} & \text{A11} & \text{Sitting down} \\
        \hline
        \textbf{C5} &  \text{Pick a pen from the
ground} & \text{A12} & \text{Pick a pen from the ground} \\
     \hline        
    \end{tabular}
    \end{adjustbox}
\end{table}


\begin{table}[H]
\scriptsize
\centering
\caption{Data files naming convention}
\label{tab:Table2}
\begin{adjustbox}{width=1\textwidth}
    \begin{tabular}{ l l l   }
    \hline
        \textbf{Symbol} & \textbf{Abbreviation for} & \textbf{Range} \\
        \hline
        \textbf{E} &  \text{Environment} & \text{\{1, 2, 3\}}   \\
        \textbf{S} &  \text{Subject} & \text{\{1, 2, 3 …, 30\}}   \\
        \textbf{C} &  \text{Experiment Class} & \text{\{1, 2, \colorbox{yellow}3, 4, 5\}}   \\
        \textbf{A} &  \text{Activity} & \text{\{1, 2, 3, …, 12\}}   \\
        \textbf{T} &  \text{Trial} & \text{\{1, 2, 3, …, 20\}}   \\
        
     \hline        
    \end{tabular}
    \end{adjustbox}
\end{table}

\subsection{Participant Details}
A total of 30 participants contributed to the dataset, offering diversity in:
\begin{itemize}
    \item Gait styles
    \item Body structure (height, weight)
    \item Motion characteristics
\end{itemize}

Each subject performed all 12 activities multiple times, providing sufficient training and testing samples for deep learning–based biometric identification.

\begin{table}[H]
\scriptsize
\centering
\caption{Details of the utilized environments along with the subject’s gender, age, weight, and height.}
\label{tab:Table3}
\begin{adjustbox}{width=1\textwidth}
    \begin{tabular}{ l l l  l l l l  }
    \hline
        \textbf{Environment ID} & \textbf{Environment Configuration} & \textbf{Subject ID} & \textbf{Gender} & \textbf{Age (years)} & \textbf{Weight (kg)} & \textbf{Height (cm)}\\
        \hline
        \textbf{1} &  \text{LOS} & \text{S1}  & \text{M}  & \text{35} & \text{75} & \text{176}\\
        \textbf{} &  \text{ } & \text{S2}  & \text{M}  & \text{21} & \text{63} & \text{188}\\
        \textbf{} &  \text{ } & \text{S3}  & \text{M}  & \text{25} & \text{75} & \text{184}\\
        \textbf{} &  \text{ } & \text{S4}  & \text{M}  & \text{23} & \text{70} & \text{178}\\
        \textbf{} &  \text{ } & \text{S5}  & \text{M}  & \text{21} & \text{77} & \text{176}\\
        \textbf{} &  \text{ } & \text{S6}  & \text{M}  & \text{20} & \text{85} & \text{180}\\
        \textbf{} &  \text{ } & \text{S7}  & \text{M}  & \text{21} & \text{90} & \text{185}\\
        \textbf{} &  \text{ } & \text{S8}  & \text{M}  & \text{21} & \text{115} & \text{187}\\
         \textbf{} &  \text{ } & \text{S9}  & \text{M}  & \text{21} & \text{90} & \text{174}\\
         \hline
        \textbf{2} &  \text{LOS} & \text{S11}  & \text{M}  & \text{23} & \text{93} & \text{180}\\
        \textbf{} &  \text{ } & \text{S12}  & \text{M}  & \text{26} & \text{60} & \text{178}\\
        \textbf{} &  \text{ } & \text{S13}  & \text{M}  & \text{20} & \text{98} & \text{188}\\
        \textbf{} &  \text{ } & \text{S14}  & \text{M}  & \text{21} & \text{104} & \text{172}\\
        \textbf{} &  \text{ } & \text{S15}  & \text{M}  & \text{25} & \text{79} & \text{175}\\
        \textbf{} &  \text{ } & \text{S16}  & \text{M}  & \text{21} & \text{70} & \text{182}\\
        \textbf{} &  \text{ } & \text{S17}  & \text{M}  & \text{24} & \text{130} & \text{185}\\
        \textbf{} &  \text{ } & \text{S18}  & \text{F}  & \text{24} & \text{60} & \text{160}\\
         \textbf{} &  \text{ } & \text{S19}  & \text{M}  & \text{23} & \text{85} & \text{175}\\
         \textbf{} &  \text{ } & \text{S20}  & \text{F}  & \text{24} & \text{62} & \text{153}\\
        \hline
        \textbf{1} &  \text{NLOS} & \text{S21}  & \text{M}  & \text{21} & \text{70} & \text{180}\\
        \textbf{} &  \text{ } & \text{S22}  & \text{M}  & \text{22} & \text{94} & \text{186}\\
        \textbf{} &  \text{ } & \text{S23}  & \text{M}  & \text{21} & \text{91} & \text{192}\\
        \textbf{} &  \text{ } & \text{S24}  & \text{M}  & \text{21} & \text{57} & \text{173}\\
        \textbf{} &  \text{ } & \text{S25}  & \text{M}  & \text{21} & \text{118} & \text{167}\\
        \textbf{} &  \text{ } & \text{S26}  & \text{M}  & \text{21} & \text{73} & \text{178}\\
        \textbf{} &  \text{ } & \text{S27}  & \text{M}  & \text{21} & \text{56} & \text{173}\\
        \textbf{} &  \text{ } & \text{S28}  & \text{M}  & \text{23} & \text{72} & \text{171}\\
         \textbf{} &  \text{ } & \text{S29}  & \text{M}  & \text{23} & \text{92} & \text{188}\\
         \textbf{} &  \text{ } & \text{S30}  & \text{M}  & \text{21} & \text{72} & \text{187}\\
        
     \hline        
    \end{tabular}
    \end{adjustbox}
\end{table}


\subsection{Data Volume Summary}
Due to repeated activity sessions across all subjects, the dataset includes:
\begin{itemize}
    \item Hundreds of CSI recordings
    \item Millions of CSI packets
    \item Several gigabytes of raw data
\end{itemize}

This volume supports the training of deep learning architectures that require large-scale time-series input.

\subsection{Suitability for This Research}
The dataset is highly suitable for CSI-based HAR and Human-ID tasks because:
\begin{itemize}
    \item CSI captures micro-movement signatures unique to each subject.
    \item LOS and NLOS conditions enable realistic performance evaluation.
    \item The dataset structure supports windowing and multivariate time-series modeling.
    \item The diversity of activities promotes generalization across motion patterns.
\end{itemize}

Thus, this dataset facilitates a rigorous comparison of multiple deep learning architectures for biometric identification and activity classification using Wi-Fi CSI.


